Gerum ráð fyrir að til sé punktur $P$ sem liggur innan allra sex hringjanna.
Látum miðjur þeirra vera $O_1, O_2, \dots, O_6$ og geisla $r_1, r_2, \dots, r_6$.
Skilyrðið að engin hringur inniheldur miðju annars þýðir að
\[
O_iO_j > \max(r_i, r_j)
\quad\text{fyrir öll } i\neq j,
\]
þannig að hver miðja liggur utan hinna hringjanna.

Þar sem $P$ er innan allra hringjanna gildir
\[
O_iP<r_i \quad \text{fyrir öll } i.
\]
Línurnar $O_1P, O_2P, \dots, O_6P$ geisla út frá sama punkti $P$.
Látum $\theta_i = \angle O_i P O_{i+1}$.

Í þríhyrningnum $O_i P O_{i+1}$ höfum við:
\[
O_iO_{i+1}\;>\;\max\{r_i,r_{i+1}\},
\quad
O_iP<r_i,\quad \text{og} \quad O_{i+1}P<r_{i+1}
\]
Því fæst
\[
O_iO_{i+1} \;>\; \max\{\,O_iP,\;O_{i+1}P\,\},
\]
svo $O_iO_{i+1}$ er stærsta hlið þríhyrningsins $O_i P O_{i+1}$.
Þar með er hornið við $P$ stærst: $\theta_i>\angle O_{i+1}O_iP$ og $\theta_i>\angle PO_{i+1}O_i$.
Setjum $\alpha_i=\angle O_{i+1}O_iP$ og $\beta_i=\angle PO_{i+1}O_i$.
Þá gildir $\alpha_i+\beta_i=180^\circ-\theta_i$ og $\theta_i>\alpha_i,\theta_i>\beta_i$, svo
\[
\theta_i \;>\; \frac{\alpha_i+\beta_i}{2} \;=\; \frac{180^\circ-\theta_i}{2}
\;\;\Rightarrow\;\; 3\theta_i>180^\circ
\;\;\Rightarrow\;\; \theta_i>60^\circ.
\]

Þar af leiðir fyrir öll sex pörin:
\[
\theta_1+\theta_2+\cdots+\theta_6 \;>\; 6\cdot 60^\circ \;=\; 360^\circ,
\]
en horn um sama punkt í planinu getur í mesta lagi numið $360^\circ$ sem gefur mótsögn.

Því er ekki til punktur $P$ sem liggur inni í öllum sex hringjunum.